\documentclass[11pt]{article}

\usepackage[T1]{fontenc}
\usepackage{amsmath,amssymb,amsthm}
\usepackage[margin=1in]{geometry}
\usepackage{microtype}
\usepackage[hidelinks]{hyperref}

\newtheorem{theorem}{Theorem}
\newtheorem{lemma}[theorem]{Lemma}

\newcommand{\Z}{\mathbb Z}
\newcommand{\ind}{\mathbf 1}

\title{Nonexistence of a Cyclic Incidence Ordering\\
for the Complete $4$-Uniform Hypergraph on Eight Vertices}
\author{Xiaojian Zeng\\
PBC School of Finance, Tsinghua University, 100190\\
\texttt{xiaojianzeng@foxmail.com}}
\date{}

\begin{document}
\maketitle

\begin{abstract}
Let the seventy $4$-subsets of an eight-element set be placed in a cyclic
order.  We prove that their eight binary incidence columns cannot all be
cyclic shifts of one common column, even when the shift of each column is
chosen independently and no adjacency condition is imposed on consecutive
$4$-sets.  In the notation of Zhou, this says that $\mathsf E(8,4)$ is false.
The coordinate-count identity from Zhou's work first confines the eight
shifts to a translate of the subgroup $7\Z_{70}$.  Seven residue slices then
give height-four factorizations on $\Z_{10}$.  Passing to the two-point
complement turns each factorization into a onefold translational tiling, whose
elementary classification permits at most sixteen distinct attendee sets,
rather than the required seventy.  A complete exact enumeration of
all $11{,}340$ reduced slice pairs is included as a reproducibility check but
is not used in place of the proof.
\end{abstract}

\section{Context and result}

Single-track Gray codes arrange binary words so that the coordinate columns
are cyclic shifts of one track, together with a prescribed adjacency rule
between consecutive words; this terminology goes back to Hiltgen, Paterson,
and Brandestini~\cite{hiltgen1996}.  For fixed-weight words, Gregor, Merino,
and M\"utze proved that if $\gcd(n,k)=1$, then the Johnson graph $J_{n,k}$ has
a balanced one-track Hamilton cycle~\cite[Theorem~4.5]{gregor2024}.  In
particular, coprimality supplies an ordering even with the extra requirement
that consecutive $k$-sets differ by one exchange.

Zhou~\cite{zhou2026} isolated the weaker, adjacency-free property in which
the shifts are independent.  Write $N=\binom nk$ and let all indices lie in
$\Z_N$.  The property $\mathsf E(n,k)$ is the existence of a bijection
\[
 e:\Z_N\longrightarrow \binom{[n]}k,
\]
a set $D\subseteq\Z_N$, and shifts $s_v\in\Z_N$ such that
\[
 \{t\in\Z_N:v\in e(t)\}=D+s_v\qquad(v\in[n]).             \tag{1}
\]
Zhou proved the necessity of coprimality for uniformities two and three and
formulated the general coprimality criterion as a conjecture.  The pair
$(n,k)=(8,4)$ is the first middle-layer instance not decided by those
classifications and complement symmetry.  We settle exactly this instance.

\begin{theorem}\label{thm:main}
There are no bijection
$e:\Z_{70}\to\binom{[8]}4$, set $D\subseteq\Z_{70}$ of size $35$, and
shifts $s_v\in\Z_{70}$ satisfying \emph{(1)}.  Equivalently,
$\mathsf E(8,4)$ is false.
\end{theorem}

The comparison above is limited to the cited sources and the literature
screened there; it is not a general priority claim.  The proof below adds to
Zhou's torsion-coset reduction a complete classification of the resulting
two-point tilings.  Its quotient-and-slice organization is also related in
spirit to coset and fiber constructions used for a different finite-abelian
problem by Zeng~\cite{zeng2026rainbow}; no result from that paper is needed
here.

\section{The coordinate-count identity}

For functions on a finite abelian group $G$, use
\[
 (F*H)(t)=\sum_{y\in G}F(y)H(t-y),
\]
and write $F^-(x)=F(-x)$.  We first record the part of
\cite[Lemma~2.2]{zhou2026} used in the proof, including its short argument so
that the reduction is self-contained.

\begin{lemma}[reflected balance]\label{lem:reflection}
If $F,H:G\to\mathbb R$ satisfy $F*H=r\ind_G$, then
$F^-*H=F*H^-=r\ind_G$.
\end{lemma}

\begin{proof}
At every nontrivial character $\chi$ of $G$,
$\widehat F(\chi)\widehat H(\chi)=0$.  Since $F$ and $H$ are real-valued,
reflection complex-conjugates their Fourier transforms and therefore
preserves their zero sets.  At the trivial character the masses are
unchanged.  Fourier inversion proves both identities.
\end{proof}

\begin{lemma}[coordinate count]\label{lem:coordinate}
Let $S,D\subseteq G$, with $|S|=m$, satisfy
$\ind_S*\ind_D=r\ind_G$.  Suppose also that
\[
 |(D+b)\cap(D+c)|=\lambda
 \qquad\text{for distinct }b,c\in S.
\]
If $T=\sum_{c\in S}c$ in $G$, then
\[
 \lambda(T-mb)=0\qquad(b\in S).                         \tag{2}
\]
\end{lemma}

\begin{proof}
Fix $b\in S$ and index all pairs $(c,t)$ with $c\in S\setminus\{b\}$
and $t\in(D+b)\cap(D+c)$.  Attach the two coordinates
$x=t-b\in D$ and $y=t-c\in D$.  For a fixed $x\in D$, balance at $b+x$
shows that $x$ occurs exactly $r-1$ times: the omitted representation is
the one with $c=b$.  For a fixed $y\in D$, the identity
$\ind_S*\ind_{-D}=r\ind_G$ from Lemma~\ref{lem:reflection}, evaluated at
$b-y$, gives the same multiplicity $r-1$ for $y$.

Thus the sum of $x-y$ over all indexed pairs is zero.  On the other hand,
$x-y=c-b$, and each $c\ne b$ indexes $\lambda$ pairs.  Hence
\[
 0=\lambda\sum_{c\in S\setminus\{b\}}(c-b)
  =\lambda(T-mb),
\]
which is (2).
\end{proof}

\section{Proof of Theorem~\ref{thm:main}}

Suppose that data as in the theorem exist.  Two vertices cannot have the
same shift: their incidence supports would then agree, whereas some
$4$-subset contains either one and excludes the other.  Consequently
\[
 S=\{s_v:v\in[8]\}\subseteq\Z_{70}
\]
has eight elements.  At each time exactly four vertices occur, so
\[
 \ind_S*\ind_D=4\ind_{\Z_{70}}.                          \tag{3}
\]
For two distinct shifts, the corresponding supports meet once for every
$4$-set containing the associated pair of vertices.  Therefore
\[
 |(D+b)\cap(D+c)|=\binom62=15\qquad(b\ne c\text{ in }S). \tag{4}
\]
Lemma~\ref{lem:coordinate}, with $m=8$, gives
\[
 15(T-8b)=0\qquad(b\in S).
\]
Subtracting the relations for $b,b'\in S$ yields
\[
 120(b-b')=0\pmod {70}.                                  \tag{5}
\]
The kernel of multiplication by $120$ on $\Z_{70}$ is exactly
$7\Z_{70}$: it has $\gcd(120,70)=10$ elements, namely the multiples of
$7$.  In particular, no nonunit has been cancelled in (5).

Fix $b_0\in S$.  Replacing $S$ by $S-b_0$ and $D$ by $D+b_0$ preserves
each individual support because
\[
 (D+b_0)+(s-b_0)=D+s.
\]
After this support-preserving normalization, (5) lets us write
\[
 S=7A,\qquad A\subseteq\Z_{10},\qquad |A|=8.             \tag{6}
\]
Relabel the vertices by the corresponding elements of $A$.

Every time has a unique expression $j+7x$, with $0\le j<7$ and
$x\in\Z_{10}$.  Define
\[
 D_j=\{x\in\Z_{10}:j+7x\in D\}.
\]
At time $j+7x$, the vertex labelled $a\in A$ occurs precisely when
\[
 j+7x\in D+7a
 \quad\Longleftrightarrow\quad x-a\in D_j.
\]
Thus its four-element attendee set is
\[
 C_{j,x}=A\cap(x-D_j),                                   \tag{7}
\]
and (3), restricted to the slice $j$, becomes
\[
 \ind_A*\ind_{D_j}=4\ind_{\Z_{10}}.                     \tag{8}
\]
Summing (8) over $x$ gives $8|D_j|=40$, so
\[
 |D_j|=5\qquad(0\le j<7).                               \tag{9}
\]

Let $B=\Z_{10}\setminus A$.  Then $B$ has two elements.  Since
$\ind_{\Z_{10}}*\ind_{D_j}=5\ind_{\Z_{10}}$, subtracting (8) gives the
onefold tiling equation
\[
 \ind_B*\ind_{D_j}=\ind_{\Z_{10}}.                     \tag{10}
\]
This is a particularly small finite-cyclic instance of a translational
tiling equation; for broader structural results on integer tiles, see
Coven and Meyerowitz~\cite{coven1999}.

Write $B=\{b,b+d\}$ with $d\ne0$, and let $f=\ind_{D_j}$.  After replacing
$x$ by $z+b$, equation (10) reads
\[
 f(z)+f(z-d)=1\qquad(z\in\Z_{10}).                       \tag{11}
\]
The values of $f$ must alternate around every orbit of translation by
$d$.  Such an alternation closes only when the order of $d$ is even.  In
$\Z_{10}$ this happens exactly when $d$ is odd.

If $d\ne5$ is odd, then $d$ generates $\Z_{10}$.  Equation (11) has only
the two parity classes as solutions for $D_j$.  The two points of $B$ have
opposite parity, so $A$ contains four elements of each parity.  For every
$x$, the translate-reflection $x-D_j$ is a parity class; hence (7) is one
of only two fixed $4$-subsets of $A$.  Already the ten times in one slice
therefore repeat attendee sets, contradicting injectivity of $e$.

It remains to consider $d=5$.  Put $H=\{0,5\}$.  The complement $B$ is one
$H$-coset and $A$ is the union of the other four.  Equation (11) says that
each $D_j$ chooses one point from every one of the five $H$-cosets.  The
same is true of every $x-D_j$.  Consequently (7) chooses one point from
each of the four fixed $H$-cosets comprising $A$.  Across all seven slices
there are therefore at most
\[
 2^4=16
\]
different attendee sets.  A bijection onto $\binom A4$ would require all
$\binom84=70$ attendee sets.  This is the final contradiction.  Even
differences make (11) impossible, and the two odd-difference cases above
are exhaustive.  The theorem follows.

\section{Exact reproducibility check}

The accompanying script \texttt{verify\_slices.py} separately enumerates
all
\[
 \binom{10}{8}\binom{10}{5}=11{,}340
\]
pairs $(A,D_j)$.  For each pair it checks every value of the convolution in
(8) and reconstructs every set in (7), using only finite sets and integer
arithmetic.  It finds twenty complements with no valid template, twenty
non-antipodal odd-difference complements with two templates and two accessible
attendee sets each, and five antipodal complements with thirty-two templates
and sixteen accessible attendee sets each.  Thus there are exactly $200$
valid reduced slice pairs, and no fixed $A$ can access all seventy
$4$-subsets.  The deterministic check can be repeated with
\begin{verbatim}
python3 verify_slices.py
\end{verbatim}
There is no random sampling or floating-point calculation.  This enumeration
corroborates the two cases in the proof; the quantified theorem follows from
the symbolic argument above.

\paragraph{Computational and AI assistance.}
OpenAI language-model tools, including Codex, were used extensively in proof
development, checking intermediate arguments, locating and comparing
literature, preparing the exact enumeration, and drafting the exposition.
The author reviewed the mathematical argument and citations and takes
responsibility for the final text.  The exact Python enumeration was used
only as the corroborating check described above.

\small
\bibliographystyle{plain}
\bibliography{references}

\end{document}
